\section{Introduction}\label{intro:section}

Random planar maps give discrete models of random surfaces. A central
problem is to identify their scaling limits in prescribed conformal
coordinates, together with area, intrinsic distance and a statistical
mechanical decoration. The continuum area measure of Liouville quantum
gravity (LQG) is a Gaussian multiplicative chaos measure. Its construction
belongs to Kahane's theory \citep{KahaneChaos1985}; Duplantier and
Sheffield developed the circle-average construction for LQG and proved
the probabilistic KPZ relation \citep{DuplantierSheffieldKPZ2011}.
For maps decorated by the critical Fortuin--Kasteleyn (FK)
random-cluster model \citep{FortuinKasteleyn1972}, the interfaces are
predicted to converge to a conformal loop ensemble (CLE), independent
of the LQG surface in its conformal coordinate.

The discrete route to this prediction begins with tree-decorated maps.
Mullin's correspondence and Bernardi's bijective developments encode
such maps through pairs of tree contours
\citep{Mullin1967,Bernardi2007}. Sheffield's inventory bijection combines
this description with Bernardi's subgraph-to-tree construction
\citep{Bernardi2008Subgraphs,SheffieldInventory}. It encodes an
FK-decorated map by a word whose two resolved contours have a correlated
Brownian scaling limit. Duplantier--Miller--Sheffield identify the
continuum mating with an LQG surface carrying an independent
space-filling Schramm--Loewner evolution (SLE); Miller--Sheffield give
the finite-area sphere version \citep{DMSMating,MillerSheffieldSphere}.
Gwynne, Mao and Sun, and Gwynne and Sun, strengthen the discrete
convergence to retain cone times and encoded FK-loop information,
including quantum areas and boundary lengths. Their finite-volume
result applies to words conditioned to reduce to the empty word
\citep{GwynneMaoSunFKI,GwynneSunFKII,GwynneSunFKIII}.

These results identify the limit through its encoding. Convergence
in a prescribed discrete conformal embedding requires a further
comparison with the embedded geometry. An intrinsic metric-space
limit alone likewise does not specify where the discrete vertices
lie in that coordinate \citep{MillerSheffieldSphere,GHS-survey}.
Gwynne, Miller and Sheffield proved convergence of the Tutte embedding
of mated-CRT maps, jointly with quantum area, the space-filling curve
and embedded random walk
\citep[Theorem~1.1]{GwynneMillerSheffieldTutte}. Holden and Sun proved
joint metric and measure convergence under the Cardy embedding for
uniform triangulations of polygons at $\gamma=\sqrt{8/3}$, together
with convergence of independent percolation decorations
\citep[Theorems~1.3 and~1.6]{HS-Cardy}. Holden and Yu proved large-scale
comparison results for circle packing and Riemann uniformization of
infinite planar maps in ergodic scale-free environments under geometric
moment and connectivity hypotheses
\citep[Theorems~1.6--1.9]{HoldenYuUniformization}. These results concern
different map laws or embedding inputs from the finite spherical FK
model studied here. A general conformal-embedding universality
conjecture is recorded as Conjecture~3.5 of \citet{HS-survey}.

Here we prove the joint statement for the flag-triangle uniformization of
critical spherical FK maps, for every fixed $0<q<4$. This establishes the
canonical conformal, area, intrinsic-metric and interface limit for this
specified model and embedding within the broader universality prediction.
Here ``subcritical'' refers to the LQG range $\gamma<2$, equivalently
$q<4$; the FK model itself is at its self-dual critical point.
The intrinsic step uses the distinct metric companion
\citep[Theorem~1.1]{Intrinsic-input}; we restate precisely its required
unmarked Gromov--Hausdorff consequence in \cref{intro:intrinsic-input}.
The new comparison identifies all these objects in the same canonical
coordinate, including atypical vertices and the order along every
macroscopic loop.

\subsection{The result}

Fix $q\in(0,4)$. Let $\gamma\in(\sqrt2,2)$ and $\kappa\in(4,8)$ be
determined by
\begin{equation}\label{intro:parameters}
 q=2+2\cos\!\left(\frac{\pi\gamma^2}{2}\right),
 \qquad \kappa=\frac{16}{\gamma^2}.
\end{equation}
The finite law and all topologies are specified in
\cref{intro:model-section}. In brief, a rooted $n$-edge decorated map
$(M_n,A_n)$ has weight $q^{\ell(M_n,A_n)/2}$, where $\ell$ counts FK
interfaces. Its $4n$ equilateral flag triangles form a conformal sphere
$X_n$ with normalized face area $m_n$. Three independent $m_n$ points
determine the orientation-preserving uniformization
$\phi_n:X_n\to\Chat$ by their images $0,1,\infty$.
Distances $d_n$ use every original primal edge. Let $\Gamma_n$ be the
images of the prescribed polygonal interfaces, with all nested loops
retained, and let
\begin{equation}\label{intro:distance-graph}
 K_n(a)=\{(\phi_n(u),\phi_n(v),a d_n(u,v)):u,v\in V(M_n)\}.
\end{equation}

On the continuum side, take the ordinary unit-area $\gamma$-LQG sphere,
forget its original marks, and normalize its conformal embedding using
three independent quantum-area points at $0,1,\infty$. Denote the
resulting area measure and intrinsic metric by $\mu_h$ and $D_h$.
Fix the deterministic normalization of $D_h$ throughout. Conditional on
this embedded surface, let $\Gamma$ have the M\"obius-invariant whole-plane
nested $\mathrm{CLE}_{\kappa}$ law on $\Chat$.

\begin{theorem}[Joint canonical limit]\label{intro:main}
There is a deterministic sequence
$a_n=a_n(q)>0$ with $a_n\to0$ such that, through all positive integers,
\begin{equation}\label{intro:joint-convergence}
 \big((\phi_n)_*m_n,K_n(a_n),\Gamma_n\big)
 \ \Longrightarrow\
 \big(\mu_h,\{(z,w,D_h(z,w)):z,w\in\Chat\},\Gamma\big).
\end{equation}
The three topologies are weak convergence of probability measures,
Hausdorff convergence of the full distance graphs, and uniform
reparameterization matching of loop collections as defined in
\cref{intro:loop-topology}. Moreover,
\begin{equation}\label{intro:mesh-conclusion}
 \max_{T\text{ a flag triangle of }X_n}
       \diam_{\rho}\phi_n(T)\ \longrightarrow\ 0
 \qquad\text{in probability},
\end{equation}
where $\rho$ is spherical distance. The conditional law of $\Gamma$ is
independent of the embedded marked quantum surface.
\end{theorem}

On a coupling with the asserted convergence, the distance-graph conclusion
controls every pair of vertices whose embedded positions converge,
including vertices with atypical degrees
or near the exploration root. The loop topology retains the order along
each curve and its multiplicity, as well as its trace. These conclusions
are simultaneous in the prescribed coordinate. All assertions are for
a fixed $q$; the proof gives no endpoint assertion, convergence rate
or explicit formula for $a_n$.

\subsection{The intrinsic input}

\begin{theorem}[Intrinsic spherical scaling; companion consequence]
\label{intro:intrinsic-input}
For the finite law of \cref{intro:finite-law} and $\gamma$ in
\eqref{intro:parameters}, there is a deterministic sequence $a_n>0$
with $a_n\to0$ such that
\[
 (V(M_n),a_n d_n)\ \Longrightarrow\ (S,D_h)
\]
in the Gromov--Hausdorff topology on compact metric spaces up to
isometry, through all positive integers. The limit is the ordinary
unit-area $\gamma$-LQG sphere of \citet[Definition~4.21(ii)]{DMSMating},
disintegrated at total quantum area one, with its marked points
forgotten and with the fixed intrinsic metric normalization. The root
and FK decoration are forgotten in this convergence.
\end{theorem}

\begin{proof}
Apply the fixed-$q$ FK clause of \citet[Theorem~1.1]{Intrinsic-input} and
forget the degree measure from its Gromov--Hausdorff--Prokhorov conclusion.
The finite law, all-edge distance, ordinary unit-area sphere law and fixed
metric normalization are exactly those above.
\end{proof}

The companion \emph{Metric-measure limits of subcritical FK and
spanning-tree planar maps}, by OpenAI, September 24, 2026, proves the
intrinsic result. We use only the displayed GH consequence; its stronger
measured conclusion and its separate spanning-tree clause are not needed
here. We take the same deterministic
$a_n(q)$ throughout. In the companion the vertex measure is
$\deg(v)/(2n)$, with loops counted twice. Each vertex belongs to
$2\deg(v)$ of the $4n$ flag occurrences, so projection of uniform flag
area to its vertex has exactly that degree distribution. The Brownian
encoding uses $2n$ letters and an area clock of duration one; neither this
clock normalization nor the flag geometry changes the all-edge metric.
The role of \cref{intro:intrinsic-input} is to complete locally identified
passage metrics at every vertex and to fix the deterministic scale.

\subsection{Proof strategy}

We first place the discrete and continuum objects on one topological
sphere. The contour limit and its continuum sewing are recalled in
\cref{enc:section}. In \cref{top:section} they give homeomorphisms
$H_n$ from the complete flag surface to the reference LQG sphere,
with every flag image small. These homeomorphisms preserve separation
and intersections but do not yet compare conformal moduli or lengths.

Here is the local information needed for those comparisons. Take a
compact disk inside a larger open region of the reference sphere.
Retain every exploration-time visit to the disk, together with a
buffer around those visits and every incidence at their primal and
dual vertices. There are finitely many retained time blocks. Their
raw words, relative contour heights and exact minimum tests determine
the corresponding part of the discrete flag surface. One can then
read, for example, the minimum graph cost between two small endpoint
sets, or the extremal length of curves joining two sides of a
quadrilateral. \Cref{local:section} constructs these finite records;
retaining all visits is what makes their conclusions apply to every
local path and vertex.

The continuum contours need not determine a subsequential discrete
readout. \Cref{ker:section} instead identifies its conditional law:
given the full continuum surface and exploration, the readout depends
only on the retained local input, and disjoint protected regions have
product conditional laws. An exact word-resampling identity shows
that changes from the bilateral word to the finite conditioned word
alter only the density of the continuum input. The same kernels then
transfer to ordinary Gaussian fields and to the sphere at its actual
area-normalized height. \Cref{ann:section} turns local events with
high probability into successful annuli at many radii around every
point. This supplies the common probabilistic input for the conformal
and metric comparisons.

\paragraph{Conformal coordinate and area.}
\Cref{conf:section} rules out degeneration of flag-surface extremal
length. A hypothetically small crossing modulus gives a probability
flow of paths with low traffic energy. Annular lower bounds force its
projected limit to consist of nonconstant rectifiable curves whose
average arclength is absolutely continuous with respect to planar area.
Locality and rotation covariance then give cheap crossings in two
transverse directions at the same point, contradicting quadrilateral
reciprocity. The resulting annular bounds give a quasiconformal limit
of $\phi_n\circ H_n^{-1}$. Its infinitesimal conformal structure is
read from local moduli; triviality of the field germ and rotation
covariance make it round. Three area marks fix the resulting conformal
map, and the area and maximum-mesh conclusions follow in that coordinate.
The analytic framework comes from Ahlfors--Beurling extremal length
\citep{AhlforsBeurling1952ConformalInvariants} and the ring estimates of
Ivrii--Markovi\'c \citep{IvriiMarkovic2019}. The elimination of
deterministic anisotropy by rotation covariance has a related use in
\citet[Section~4.4]{HoldenYuUniformization}, where determinism comes
from averaged Dirichlet energies.

\paragraph{Ordered nested interfaces.}
\Cref{loop:section} identifies the loops independently of the metric
comparison. Flexible-order matches form a noncrossing chord system
on the exploration-time circle. Its faces encode the finite FK loops,
including their cyclic traversal order. Chord convergence first gives
ordered limits of these face contours. The alternating CLE exploration
identifies their initial portions; winding tests on both sides exclude
an omitted moving portion. This matches every macroscopic loop, with
multiplicity, and retains the full nested collection. The exploration
convention follows \citet{SheffieldExplorationTrees,MillerSheffieldIGIV}
and the nested construction described in
\citet[Section~2.1.6]{AruHoldenPowellSunCritical}.

\paragraph{Local passages and all-vertex distances.}
In \cref{met:section}, graph passages between small endpoint sets
first acquire upper and lower comparisons with the reference LQG
metric. Repeated annular opportunities and local finite-energy field
changes bring a reference geodesic near a shortcut. A weighted
change-of-measure argument forces the optimal comparison constants
to agree. This adapts the optimal-constant method of Gwynne and
Miller \citep[Section~1.5 and Sections~3--6]{GwynneMillerMetric} to
the local passage kernels. The changed laws are compared before a
shortcut or a geodesic hit is selected, so restriction to a rare
selected event preserves the required domination.

These passages have freely chosen endpoints in small sets; reaching
a prescribed lattice vertex is a separate problem.
\Cref{comp:section} uses \cref{intro:intrinsic-input} at this point.
Small surrounding circuits and the sphere topology of the GH limit
exclude nontrivial spatial fibers. A possible collapse of all local
passage costs would concentrate the remaining metric at one spatial
point; the joint local laws exclude it by placing two separated
positive escape-cost tests in the same sphere. The resulting
all-vertex metric is a deterministic multiple of $D_h$, and the
intrinsic diameter law fixes that multiple. All field-law arguments
use the original oriented reference coordinate; only then are their
geometric conclusions transported to the canonical coordinate.
\Cref{joint:section} assembles the three branches, including their
common marks and possible common reflection.
Figure~\ref{intro:roadmap-figure} records these dependencies.

\begin{figure}[ht]
\centering
\begin{tikzpicture}[
 box/.style={draw=linkblue,rounded corners=2pt,align=center,
 text width=2.8cm,minimum height=1cm,font=\small},
 arr/.style={-{Stealth[length=2mm]},thick,draw=linkblue},
 node distance=7mm and 9mm]
\node[box,text width=8.4cm] (common) at (0,0)
 {Contour topology and protected local data\\Conditional laws and aligned annuli};
\node[box] (conf) at (-2.8,-1.8) {Conformal coordinate,\\area and mesh};
\node[box] (pass) at (2.8,-1.8) {Local metric\\passages};
\node[box] (loops) at (-2.8,-3.6) {Ordered nested\\loop collection};
\node[box] (comp) at (2.8,-3.6) {All-vertex metric\\completion};
\node[box,text width=2.2cm] (gh) at (6.2,-3.6) {Intrinsic GH\\input};
\node[box] (joint) at (0,-5.4) {Joint canonical\\limit};
\draw[arr] (common.south west)--(conf.north);
\draw[arr] (common.south east)--(pass.north);
\draw[arr] (common.west)--(-5,0)--(-5,-3.6)--(loops.west);
\draw[arr] (conf)--(loops);
\draw[arr] (conf.south east)--(comp.north west);
\draw[arr] (loops)--(joint);
\draw[arr] (pass)--(comp);
\draw[arr] (gh)--(comp);
\draw[arr] (comp)--(joint);
\end{tikzpicture}
\caption{The conformal, loop and metric arguments share the contour
correspondence. The intrinsic GH input enters only the all-vertex
completion of the metric.}
\label{intro:roadmap-figure}
\end{figure}
